\documentclass[12pt, a4paper]{article}

\usepackage[margin=2.5cm]{geometry}

% Pacotes básicos 
% encoding
%--------------------------------------
\usepackage[T1]{fontenc}		% Seleção de códigos de fonte.
\usepackage[utf8]{inputenc}		% Codificação do documento (conversão automática dos acentos)
%--------------------------------------
% ----------------------------------------------------------
% \usepackage{lmodern}			% Usa a fonte Latin Modern	
\usepackage{fontspec}
\usepackage{Acorn}
\usepackage{lettrine}

% \setmainfont{QTCloisteredMonk}
% \setmainfont{QTLondonScroll}
% \setmainfont{QTHeidelbergType}
% \usepackage{Typocaps}
% \usepackage{ebgaramond}
% \usepackage{libertine}

\usepackage[dvipsnames]{xcolor}
% \renewcommand{\LettrineFontHook}{\color{BrickRed}\Typocapsfamily{}}
\renewcommand{\LettrineFontHook}{\color{BrickRed}\Acornfamily{}}
\setcounter{DefaultLines}{4}
\setmainfont{EB Garamond}[Numbers={Proportional},RawFeature={+ss06}]

% \newcommand*\initfamily{\usefont{U}{Acorn}{xl}{n}}

% Hyphenation rules
%--------------------------------------
\usepackage{hyphenat}
\hyphenation{mate-mática recu-perar}
%--------------------------------------
\usepackage{indentfirst}		% Indenta o primeiro parágrafo de cada seção.
\usepackage{color}		    	% Controle das cores
\usepackage{graphicx}			% Inclusão de gráficos
\graphicspath{ {./images/} }      % define diretório raiz para imagens, então é só colocar o nome do arquivo em \includegraphics{nomedoarquivo}
\usepackage{microtype} 			% para melhorias de justificação
\usepackage{ifthen}		    	% para montar condicionais
\usepackage[brazil]{babel}		% para utilizar termos em português
\usepackage[autostyle]{csquotes}    % fazer as aspas aparecerem corretamente
% \MakeOuterQuote{"}
% \usepackage[final]{pdfpages}    % para incluir páginas de arquivos PDF
% \usepackage{lipsum}				% para geração de lorem ipsum
\usepackage{float} 		        % Para a figura ficar na posição correta
\usepackage{caption}            % configura legenda 
    \captionsetup{font=small}	% tamanho da fonte 10pt
\usepackage{footmisc}           % precisa ir antes de hyperref, senao referencias a rodapes quebram
\usepackage{hyperref}
% \usepackage[colorlinks=true]{hyperref}


% \usepackage{enotez}      % <-- instead of \usepackage{endnotes}
% \let\footnote=\endnote

% \makeatletter
% \def\enoteheading{\section*{\notesname
%   \@mkboth{\MakeUppercase{\notesname}}{\MakeUppercase{\notesname}}}%
%   \mbox{}\par\vskip-2.3\baselineskip\noindent\rule{.5\textwidth}{0.4pt}\par\vskip\baselineskip}
% \makeatother

% \usepackage[
% backend=biber,
% style=numeric,
% sorting=none                    % sort pela ordem de ocorrência
% ]{biblatex}                     % usa o biblatex, ferramenta + moderna para a bibliografia
%     \addbibresource{references.bib}     % puxa as entries de bibliografia desse arquivo


% https://tex.stackexchange.com/questions/13756/quote-environment-with-reference-at-the-end-right
\def\signed #1{{\leavevmode\unskip\nobreak\hfil\penalty50\hskip2em
  \hbox{}\nobreak\hfil(#1)%
  \parfillskip=0pt \finalhyphendemerits=0 \endgraf}}

\newsavebox\mybox
\newenvironment{aquote}[1]
  {\savebox\mybox{#1}\begin{quote}}
  {\signed{\usebox\mybox}\end{quote}}
  
% https://tex.stackexchange.com/a/74609
\usepackage{fourier-orns}
\newcommand{\ornamento}{
\noindent
    \parbox{\textwidth}{%
        {\color{BrickRed}\raisebox{-2.1pt}[10pt][10pt]{\leafright} \hrulefill \raisebox{-2.1pt}[10pt][10pt]{~~~\decofourleft \decotwo\decofourright~~~} \hrulefill \raisebox{-2.1pt}[10pt][10pt]{ \leafleft}}
    }%
}

\usepackage{tcolorbox}
% Criação de um ambiente padrão para alertas e armadilhas comuns
\newtcolorbox{atencao}[1][Ponto de Atenção]{
  colback=red!5!white,      % Fundo levemente avermelhado
  colframe=red!75!black,    % Borda vermelha escura
  fonttitle=\bfseries,      % Título em negrito
  title=#1                  % Permite mudar o título se necessário
}

\title{Manual do Template \LaTeX \ PPGT \\ Guia de Transição e Uso}
\author{Samuel Seixas Martin}
\date{23 de outubro de 2026}

% first section starts at 0
\renewcommand{\theenumi}{\Roman{enumi}}

\begin{document}
\maketitle
% \tableofcontents
\ornamento

\section{Introdução: A Transição do Microsoft Word}

\lettrine{O}{} objetivo primário desta base de código é eliminar a dependência de editores visuais como o Microsoft Word. O PPGT possui diretrizes tipográficas rigorosas que frequentemente sofrem desconfiguração acidental em processadores de texto convencionais. O \LaTeX \ retira do autor a carga mental de microgerenciar a formatação página por página, permitindo foco total na pesquisa, na matemática e na qualidade do texto.

\subsection{Quando o Orientador Exigir Arquivo Word (.docx): A Conversão Direta}

Um receio frequente entre novos pesquisadores é: \textit{``E se meu orientador ou coautor não usar \LaTeX\ e exigir o texto em formato Microsoft Word (.docx) para fazer anotações?''}.

Existe um método nativo, simples e de altíssima fidelidade que preserva integralmente a diagramação do documento sem necessitar de conversores externos ou softwares pagos:

\begin{enumerate}
    \item Abra o aplicativo Microsoft Word no seu computador com a tela inicial vazia (isto é, com \textbf{nenhum documento aberto}).
    \item Localize o arquivo \texttt{principal.pdf} compilado no seu gerenciador de arquivos.
    \item \textbf{Arraste o arquivo PDF diretamente para o centro da janela principal do Word} e solte.
    \item O Microsoft Word fará uma conversão automática e direta, reconstruindo parágrafos, tabelas, equações e títulos em um documento \texttt{.docx} editável com máxima fidelidade visual.
\end{enumerate}

Dessa forma, você pode manter seu fluxo de trabalho profissional e rigoroso no \LaTeX, gerando uma cópia em \texttt{.docx} apenas quando estritamente requisitado por terceiros.

\section{Principais Dificuldades na Adaptação}

A migração de um processador de texto visual para um sistema de preparação de documentos algorítmico impõe desafios iniciais notórios:

\begin{itemize}
    \item \textbf{Paradigma Visual vs. Lógico:} No Word, opera-se no modelo WYSIWYG (\textit{What You See Is What You Get}), onde a formatação é aplicada diretamente na tela. No \LaTeX, trabalha-se com estrutura semântica (declarando o que é um capítulo, uma citação ou uma equação), confiando ao motor tipográfico a responsabilidade pelo layout final. O autor precisa abrir mão da ilusão de ``controle manual'' em prol da consistência sistêmica.
    \item \textbf{Curva de Aprendizado e Depuração de Erros:} A regra de que cada chave aberta \texttt{\{} precisa ter sua correspondente fechada \texttt{\}}, declarar ambientes estritos (\texttt{\textbackslash begin} e \texttt{\textbackslash end}) e compilar o código repetidamente exige uma postura metódica. Diferente do Word, onde erros tipográficos apenas sujam a página, no \LaTeX \ um erro de sintaxe pode interromper toda a compilação.
    \item \textbf{Gestão de Elementos Flutuantes (Floats):} Tabelas, quadros e figuras não ficam ``presos'' exatamente na linha de código onde foram declarados. O algoritmo do \LaTeX \ calcula matematicamente o local ideal (topo, base da página ou página isolada) para não gerar espaços em branco ou quebrar o fluxo de leitura, o que muitas vezes frustra quem está acostumado a ``empurrar'' imagens no texto com a tecla Enter (isso é muito bom também, porque sabemos como é frustrante no Word tentar manter o layout manualmente, especialmente quando adicionamos mais texto).
\end{itemize}

\section{Estruturação Lógica: Capítulos e Seções}

Uma das diferenças mais drásticas na transição do Word para o \LaTeX \ é o abandono da formatação visual em favor da estruturação semântica. No Word, o autor frequentemente seleciona o texto e aplica formatações manuais (negrito, tamanho de fonte) para simular um título. No \LaTeX, o usuário apenas declara a função lógica daquele bloco de texto, e o motor tipográfico consulta o arquivo de estilo para aplicar as normas do PPGT de forma padronizada.

A hierarquia estrutural do documento é construída utilizando a seguinte cascata de comandos:

\begin{itemize}
    \item \textbf{Nível Primário (Capítulos):} Declarado via \texttt{\textbackslash chapter\{Título\}}. Este comando inicia automaticamente uma nova página e aplica a formatação primária (ex: \textbf{1 INTRODUÇÃO}).
    \item \textbf{Nível Secundário (Seções ou Subcapítulos):} Declarado via \texttt{\textbackslash section\{Título\}}. Subdivide o capítulo recebendo o segundo nível de numeração (ex: \textbf{1.1 Objetivo}).
    \item \textbf{Nível Terciário (Subseções):} Declarado via \texttt{\textbackslash subsection\{Título\}}. Desdobra a seção primária (ex: \textbf{1.1.2 Objetivos Específicos}).
    \item \textbf{Nível Quaternário (Subsubseções):} Declarado via \texttt{\textbackslash subsubsection\{Título\}}. Representa a ramificação mais profunda recomendada no texto (ex: \textit{(a) Dominar o mundo}).
\end{itemize}

\textbf{Regras de Ouro da Estruturação:}
\begin{enumerate}
    \item Jamais digite os números manualmente (ex: \texttt{\textbackslash section\{1.1 Minha Seção\}}). O \LaTeX \ gerencia e atualiza toda a numeração dinamicamente; basta inserir o texto do título.
    \item Nunca quebre a cadeia hierárquica. Não abra uma \texttt{\textbackslash subsection} sem antes ter declarado uma \texttt{\textbackslash section}.
    \item Não digite os títulos inteiramente em CAIXA ALTA dentro do comando. Escreva normalmente (ex: \texttt{\textbackslash chapter\{Revisão Bibliográfica\}}). Se a norma exigir maiúsculas no documento final, o template converterá o texto automaticamente durante a compilação, mantendo o seu código limpo e preservando a capitalização correta para o sumário.
    \item \textbf{Apêndices e Anexos:} Os arquivos já vêm prontos no template (\texttt{zz-apendices.tex} e \\ \texttt{zz-anexos.tex}). Se quiser colocar mais de um apêndice ou anexo, não precisa criar novos arquivos: basta adicionar mais um comando \texttt{\textbackslash chapter\{Título\}} e um rótulo \texttt{\textbackslash label} dentro deles. O \LaTeX \ gera as letras sequenciais (APÊNDICE B, ANEXO B, etc.) e atualiza o Sumário automaticamente.
\end{enumerate}

\section{Sintaxe Básica e Formatação Textual}

Para os iniciantes na linguagem \LaTeX \, alguns comandos fundamentais substituem os botões tradicionais da interface gráfica:

\begin{itemize}
    \item \textbf{Comentários (\texttt{\%}):} Qualquer caractere inserido após o símbolo de porcentagem em uma linha é ignorado pelo compilador. É a ferramenta ideal para deixar anotações estruturais no código ou desativar trechos de texto temporariamente sem apagá-los.
    
    \begin{atencao}[O Terror do Símbolo de Porcentagem]
    Embora o \texttt{\%} seja extremamente útil para deixar comentários no código, ele é a causa número um de frases ``desaparecidas'' para os iniciantes. 

    Se você quiser escrever que \textit{``a taxa de acidentes cresceu 50\% neste trecho''} e esquecer de escapar o caractere (digitando apenas \texttt{50\%}), o LaTeX transformará todo o resto do seu parágrafo em um comentário invisível a partir daquele ponto. Sempre utilize a barra invertida para imprimir o símbolo textualmente: \texttt{50\textbackslash\%}.
    \end{atencao}
    
    \item \textbf{Caracteres Reservados:} Similarmente ao caractere de porcentagem (\texttt{\%}), outros símbolos como o \textit{E comercial} (\texttt{\&}), cifrão (\texttt{\$}), \textit{underline} (\texttt{\_}) e \textit{hashtag} (\texttt{\#}) possuem funções matemáticas e estruturais no motor do \LaTeX. Se você digitá-los soltos no texto (ex: ``R\$ 100,00'' ou ``P\&D''), a compilação falhará. É obrigatório ``escapá-los'' com uma barra invertida antes: \texttt{\textbackslash\&}, \texttt{\textbackslash\$}, \texttt{\textbackslash\_}, \texttt{\textbackslash\#}.

    \item \textbf{Itálico (\texttt{\textbackslash textit\{...\}}):} Inclina a fonte. Segundo as normas do PPGT, palavras em língua estrangeira à língua principal do texto devem obrigatoriamente ser escritas em itálico.
    \item \textbf{Negrito (\texttt{\textbackslash textbf\{...\}}):} Aplica peso à fonte. 
    
    \begin{atencao}[Evitar negrito]
    Atenção ao rigor do PPGT: o manual especifica que palavras que se desejam dar destaque devem ser escritas entre aspas duplas ou com as iniciais em maiúsculo, e não com o negrito, que é reservado estritamente para títulos de seções, numeração de ilustrações e elementos estruturais.
    \end{atencao}

    \item \textbf{Listas Não Ordenadas (\texttt{itemize}):} Cria listas com marcadores (bullet points). Utiliza-se o ambiente \texttt{\textbackslash begin\{itemize\}} e \texttt{\textbackslash end\{itemize\}}, declarando cada nova linha interna com o comando \texttt{\textbackslash item}.
    \item \textbf{Listas Numeradas (\texttt{enumerate}):} Cria listas sequenciais e numeradas automaticamente. O funcionamento lógico é idêntico ao ambiente de marcadores, bastando usar \texttt{\textbackslash begin\{enumerate\}} e \texttt{\textbackslash end\{enumerate\}}.
    \item \textbf{Aspas Duplas (\texttt{``...''}):} O caractere de aspas retas do teclado (\verb|"|) não funciona corretamente no LaTeX. Para criar aspas duplas tipograficamente exatas, utilize duas crases juntas para abrir (\verb|``|) e dois apóstrofos simples para fechar (\verb|''|).
    \item \textbf{Aspas Simples (\texttt{`...'}):} Para criar aspas simples tipograficamente corretas, utilize uma crase (\verb|`|) para abrir e um apóstrofo simples (\verb|'|) para fechar. O manual do PPGT recomenda o uso de aspas simples para citações dentro de citações, mantendo a consistência tipográfica.
    \item \textbf{Parágrafos e Quebras de Linha:} O \LaTeX \ separa drasticamente a forma como você digita da forma como o documento é impresso. No Word, apertar a tecla \texttt{Enter} quebra a linha imediatamente, mas se você fizer o mesmo no \LaTeX, você será solenemente ignorado e o texto continuará fluindo no mesmo parágrafo no PDF, porque uma única quebra de linha no código-fonte é para critérios de organização do código e não para espaçamento visual no documento final.

    Para iniciar um novo parágrafo e garantir que o recuo da primeira linha seja aplicado corretamente (conforme exigido pelas normas do PPGT), é estritamente necessário \textbf{deixar uma linha em branco no código} (ou seja, apertar \texttt{Enter} duas vezes). Cada bloco de texto separado por uma linha em branco no código será interpretado como um novo parágrafo autônomo.

    \begin{atencao}[A Ilusão do Espaçamento e a Quebra Forçada]
    A transição do Word geralmente traz dois vícios de formatação que quebram o documento no LaTeX:

    \textbf{1. Apertar Enter várias vezes para criar espaço:} No Word, se você apertar \texttt{Enter} dez vezes, criará um buraco enorme na página. No LaTeX, não importa se você deixou uma ou vinte linhas em branco no código; o compilador entenderá apenas que um novo parágrafo começou e aplicará exatamente o mesmo espaçamento padrão da norma. O motor tipográfico ignora espaços vazios redundantes.

    \textbf{2. O uso indevido do \texttt{\textbackslash\textbackslash}:} Evite usar duas barras invertidas (\texttt{\textbackslash\textbackslash}) no meio de um texto corrido. Elas forçam uma quebra de linha bruta, jogando o texto para baixo sem o recuo semântico do novo parágrafo. O comando \texttt{\textbackslash\textbackslash} não significa ``novo parágrafo'', mas sim ``corte esta linha aqui'', devendo ser reservado para tabelas ou equações matemáticas.

    \vspace{0.3cm}
    \textbf{O QUE FAZER ENTÃO?} \\
    Se você realmente precisar manipular o espaço vertical de forma manual, utilize os comandos estruturais corretos:
    \begin{itemize}
        \item \textbf{Para pular de página:} Use \texttt{\textbackslash newpage} ou \texttt{\textbackslash clearpage}. (Vale lembrar que o comando \texttt{\textbackslash chapter} já inicia uma nova página automaticamente).
        \item \textbf{Para adicionar um espaço específico:} Use \texttt{\textbackslash vspace\{2cm\}} para forçar exatamente 2 centímetros de espaço vertical entre dois elementos.
        \item \textbf{Para empurrar texto para o rodapé:} Use \texttt{\textbackslash vfill}. Esse comando funciona como uma ``mola'', preenchendo todo o espaço vazio disponível e empurrando o conteúdo seguinte para a margem inferior da página.
    \end{itemize}
    \end{atencao}

    \item \textbf{Fórmulas Matemáticas:} Para inserir equações na mesma linha do texto, envolva a expressão com cifrões (ex: \texttt{\$E=mc\textasciicircum 2\$}). Para equações destacadas, centralizadas e numeradas automaticamente, utilize o ambiente \texttt{\textbackslash begin\{equation\}} e \texttt{\textbackslash end\{equation\}}.
    \item \textbf{O Desafio das Tabelas:} A formatação nativa de tabelas no \LaTeX \ pode ser frustrante inicialmente. A recomendação de ouro é utilizar ferramentas visuais online, como o \href{https://www.tablesgenerator.com/}{TablesGenerator.com}. Você constrói sua tabela graficamente (ou cola do Excel) e o site gera o código \LaTeX \ pronto para ser copiado para o seu documento.

    \item \textbf{Referências Cruzadas Inteligentes (\texttt{\textbackslash autoref}):} Jamais digite textualmente nomes e números como ``Figura 1'' ou ``Tabela 2''. O \LaTeX \ atualiza a numeração dinamicamente de acordo com a posição dos elementos. Adicione um \texttt{\textbackslash label\{nome-do-marcador\}} imediatamente após criar a legenda da sua figura ou tabela, e cite-a no texto exclusivamente com \\ \texttt{\textbackslash autoref\{nome-do-marcador\}}. O comando \texttt{\textbackslash autoref} identifica automaticamente a natureza do elemento e imprime o prefixo correto com a inicial maiúscula (ex: Figura 1, Tabela 2, Seção 3.1), transformando tudo em um link clicável e evitando inconsistências no texto.
    \item \textbf{Hifenização e Quebra de Palavras:} O motor moderno do \LaTeX \ (com os pacotes \texttt{babel} em português e codificação \texttt{T1}) separa silabicamente de forma automática quase 100\% das palavras da língua portuguesa com acentuação correta. Caso surja alguma palavra estrangeira em inglês, nome de software ou termo composto que quebre incorretamente na margem direita (ou cause aviso de \textit{Overfull \textbackslash hbox}), você dispõe de duas soluções:
    \begin{itemize}
        \item \textbf{Hífen discricionário no texto (\texttt{\textbackslash-}):} Insira \texttt{\textbackslash-} exatamente onde a palavra pode ser dividida (ex.: \texttt{feed\textbackslash-back} ou \texttt{bench\textbackslash-mark}). O \LaTeX \ só exibirá o hífen se a palavra realmente precisar ser quebrada no fim da linha.
        \item \textbf{Arquivo global de hifenização (\texttt{formatacao/hifenizacao.tex}):} Para termos técnicos repetidos ao longo do trabalho, declare a regra uma única vez no arquivo \\ \texttt{formatacao/hifenizacao.tex} com o comando \texttt{\textbackslash hyphenation\{feed-back\}}. O compilador aplicará essa separação em todo o documento.
    \end{itemize}
\end{itemize}

\section{Configuração Inicial do Documento Principal}

Ao iniciar a redação da sua dissertação, o primeiro passo é preencher os metadados estruturais no preâmbulo do seu arquivo principal (geralmente \texttt{principal.tex}). O template foi arquitetado para receber essas informações essenciais uma única vez e distribuí-las automaticamente para a Capa, Folha de Rosto, Ficha Catalográfica e Folha de Aprovação.

Localize e altere os seguintes comandos de declaração com os seus dados:

\begin{itemize}
    \item \textbf{Título do Trabalho:} Utilize \texttt{\textbackslash title\{Seu Título Aqui\}}. Se houver subtítulo, a norma exige que seja separado do título principal por dois pontos (:).
    \item \textbf{Nome e Contato do Autor:} A lógica de partículas e Ficha Catalográfica exige que o nome seja dividido em \texttt{\textbackslash ultimosobrenomeautor\{...\}} e \texttt{\textbackslash prenomesautor\{...\}}. Preencha também \texttt{\textbackslash autoremail\{...\}} e \texttt{\textbackslash autorendereco\{...\}}.
    \item \textbf{Palavras-chave:} Declare os termos indexadores individualmente usando os comandos \\ \texttt{\textbackslash palavrachavei\{...\}} até \texttt{\textbackslash palavrachavevi\{...\}}.
    \item \textbf{Orientador, Coorientador e Banca Examinadora:} Preencha os dados do orientador (nome direto para a folha de rosto, invertido para a ficha catalográfica e formato completo da banca em \texttt{\textbackslash membrodabancai\{...\}}). Caso haja coorientador, basta descomentar o bloco de 4 linhas correspondente (ele é opcional e vem comentado por padrão). Os demais examinadores são configurados em \texttt{\textbackslash membrodabancaii} e \texttt{\textbackslash membrodabancaiii}.
    \item \textbf{Data da Defesa:} O template possui um extrator automático. Preencha rigorosamente no formato \texttt{\textbackslash date\{AAAA-MM-DD\}} (ex: \texttt{2026-10-02} para o dia 2 de outubro de 2026), e o sistema separará dia, mês e ano de forma automática nas páginas iniciais.
    \item \textbf{Dados Acadêmicos e Institucionais:} Defina o escopo do documento preenchendo os comandos \texttt{\textbackslash grau\{Mestre\}}, \texttt{\textbackslash area\{Transportes\}}, \texttt{\textbackslash siglaarea\{ENC\}}, \\ \texttt{\textbackslash tipodemonografia\{Dissertação\}} e \texttt{\textbackslash titulacao\{Mestrado\}}.
    \item \textbf{Número da Publicação:} O campo \texttt{\textbackslash numpublicacao\{...\}} é opcional, caso haja necessidade.
    \item \textbf{Regras de Hifenização:} O comando \texttt{\textbackslash input\{formatacao/hifenizacao\}} inclui o arquivo \texttt{formatacao/hifenizacao.tex}, onde você pode adicionar regras de separação silábica para estrangeirismos e termos técnicos incomuns à língua portuguesa.
    \item \textbf{Seleção de Idioma (Português vs. Inglês) e Resumo Estendido:} Por padrão, o documento é redigido em português. Caso você vá defender sua dissertação em inglês, basta descomentar a linha \texttt{\textbackslash selectlanguage\{english\}} no preâmbulo do \texttt{principal.tex}. O template adaptará automaticamente todos os títulos estruturais (\textit{Chapter}, \textit{Table}, \textit{Figure}, \textit{Contents}, \textit{Appendix}, etc.). Conforme Resolução do CEPE/UnB, pesquisas em língua inglesa devem obrigatoriamente conter um \textbf{Resumo Estendido em Língua Portuguesa} (2 a 5 páginas), bastando descomentar a linha \texttt{\textbackslash include\{capitulos-pretextuais/00-resumoestendido\}} na seção pré-textual (linha 170).
\end{itemize}

\section{Imagens, Citações e Nomenclaturas}

O template conta com atalhos e pacotes configurados para otimizar o fluxo de trabalho e garantir aderência às normas sem esforço manual:

\begin{itemize}
    \item \textbf{Diretório de Imagens:} A declaração \texttt{\textbackslash graphicspath\{imagens/\}} no preâmbulo define a pasta raiz para os arquivos gráficos. Assim, basta usar o nome do arquivo (ex: \\ \texttt{\textbackslash includegraphics\{foto.png\}}). 
    
    \begin{atencao}[Nomes de Arquivos de Imagem]
    A mesma regra dos arquivos \texttt{.tex} se aplica rigorosamente aqui: jamais utilize espaços, acentos ou cedilha nos nomes das suas imagens. Um arquivo chamado \texttt{grafico fluxo.png} fará a compilação falhar silenciosamente. Renomeie sempre para formatos como \texttt{grafico-fluxo.png}.
    \end{atencao}
    
    \item \textbf{Numeração Automática (O fim da neurose):} No Word, inserir uma nova figura no meio do capítulo gera a neurose de ter que caçar e renumerar manualmente todas as imagens e menções subsequentes. No \LaTeX, isso não existe. O compilador numera figuras, tabelas e quadros de forma autônoma. Basta adicionar um \texttt{\textbackslash label\{marcador\}} logo após a legenda e chamá-la no texto com \texttt{\textbackslash autoref\{marcador\}}.
    
    \item \textbf{Gestão Bibliográfica e Citações:} As suas referências não são redigidas manualmente no texto. Elas ficam armazenadas estruturalmente em um arquivo com extensão \texttt{.bib} codificado em UTF-8. Recomenda-se o uso de gerenciadores como Mendeley, Zotero ou a exportação do Google Acadêmico para popular esse arquivo. No texto, a formatação muda de acordo com a chamada:
    \begin{itemize}
        \item Use \texttt{\textbackslash citet\{chave\}} para citações \textbf{TEXTUAIS} (citação indireta integrada). O autor faz parte da frase e apenas o ano fica entre parênteses.\\ \textit{Exemplo no código:} Segundo \texttt{\textbackslash citet\{davis2004\}}, os acidentes...\\ \textit{Resultado no PDF:} Segundo Davis (2004), os acidentes...
        \item Use \texttt{\textbackslash citep\{chave\}} para citações com \textbf{PARÊNTESES} (no fim da afirmação). O autor e o ano ficam isolados entre parênteses no final da ideia.\\ \textit{Exemplo no código:} ... na modelagem de acidentes \texttt{\textbackslash citep\{davis2004\}}.\\ \textit{Resultado no PDF:} ... na modelagem de acidentes (DAVIS, 2004).
    \end{itemize}
    
    \begin{atencao}[Localização Automática do Arquivo .bib e Como Alterar]
    O arquivo de estilo \texttt{formatacao/ppgt-unb-format.sty} detecta e inclui automaticamente o seu arquivo de referências na seguinte ordem de busca prioritária:
    \begin{enumerate}
        \item \texttt{bibliografia/referencias.bib} (local padrão recomendado para manter a raiz limpa);
        \item \texttt{referencias.bib} (na raiz do projeto, mantido por compatibilidade).
    \end{enumerate}

    \textbf{Como alterar o nome ou usar múltiplos arquivos \texttt{.bib}:}\\
    Graças ao motor moderno \textbf{BibLaTeX + Biber}, você tem total flexibilidade:
    \begin{itemize}
        \item Se o seu arquivo tiver outro nome (ex.: \texttt{minhas\_referencias.bib}), basta adicionar no preâmbulo do seu \texttt{principal.tex} a instrução:
\begin{verbatim}
\addbibresource{bibliografia/minhas_referencias.bib}
\end{verbatim}
        \item Você pode inclusive dividir suas referências em arquivos temáticos (por exemplo, \texttt{artigos.bib} e \texttt{livros.bib}) adicionando múltiplas linhas \texttt{\textbackslash addbibresource\{...\}}. O Biber unificará todas as bases e construirá a lista final em estrita ordem alfabética automaticamente.
    \end{itemize}

    \textbf{Começando sua bibliografia do zero (base limpa):}\\
    As referências que acompanham o arquivo padrão \texttt{bibliografia/referencias.bib} são apenas exemplos didáticos incluídos para demonstrar a formatação da norma do PPGT. Quando for escrever a sua dissertação real, \textbf{você pode apagar todo o conteúdo daquele arquivo} e colar diretamente as referências exportadas do seu gerenciador (Zotero, Mendeley, etc.). Apenas lembre-se de remover ou substituir as citações de teste usadas no texto dos capítulos para não gerar avisos de chaves não encontradas.
    \end{atencao}
    
    \item \textbf{Siglas e Símbolos:} Para gerar as listas pré-textuais, declare os termos ao longo do texto usando o comando \texttt{\textbackslash nomenclature\{<símbolo ou sigla>\}\{<descrição>\}}. O template oferece os atalhos \texttt{\textbackslash simbolo} e \texttt{\textbackslash sigla}. Fique tranquilo: o motor do \LaTeX \ cuidará de organizar todas as entradas em estrita ordem alfabética de forma autônoma durante a compilação.
\end{itemize}

\section{Estrutura e Inclusão de Capítulos}

Para manter o código organizado e facilitar a manutenção, o texto da dissertação não é escrito diretamente no arquivo principal. O documento é modularizado, com cada capítulo isolado em seu próprio arquivo \texttt{.tex} dentro da pasta \texttt{capitulos/}.

No meio do arquivo \texttt{principal.tex}, você encontrará o bloco responsável por unir todas essas partes:

\begin{verbatim}
  \chapter{INTRODUÇÃO}
\label{introducao}
\label{capitulo1}

% NOTA NORMATIVA PPGT: O manual oficial do PPGT NÃO prevê resumo no início dos capítulos.
% Este comando (\resumodocapitulo) é um recurso opcional herdado do modelo genérico da FT-UnB.
% Caso o seu orientador aprove ou solicite uma epígrafe/síntese no início do capítulo,
% basta descomentar a linha abaixo:
% \resumodocapitulo{Este capítulo apresenta a contextualização geral da pesquisa, a delimitação do problema e as hipóteses investigadas, estabelecendo os objetivos primários e a estrutura lógica da dissertação.}

Apresente aqui a introdução geral da sua pesquisa, contextualizando o tema de estudo no âmbito dos transportes e indicando a relevância do trabalho desenvolvido.

\section{DELIMITAÇÃO DO PROBLEMA}
\label{sec:problema}

Descreva a delimitação clara do problema investigado. Esclareça quais aspectos da realidade serão abordados e quais lacunas práticas ou teóricas motivam a investigação.

\section{HIPÓTESES}
\label{sec:hipoteses}

Apresente as hipóteses de pesquisa formuladas, indicando as relações esperadas entre as variáveis analisadas e os fenômenos observados.

\section{OBJETIVOS}
\label{sec:objetivos}

\subsection{Objetivo Geral}

Apresente o objetivo geral da pesquisa de forma clara, direta e iniciada por verbo no infinitivo.

\subsection{Objetivos Específicos}

Liste as etapas e metas intermediárias necessárias para atingir o objetivo geral:
\begin{enumerate}[(a)]
  \item Primeiro objetivo específico da pesquisa;
  \item Segundo objetivo específico da pesquisa;
  \item Terceiro objetivo específico da pesquisa.
\end{enumerate}

\section{JUSTIFICATIVA}
\label{sec:justificativa}

Apresente a justificativa técnica, científica, socioeconômica e ambiental da pesquisa, evidenciando as contribuições práticas e acadêmicas esperadas para a área de Transportes.

A seguir é apresentado um exemplo de elemento do tipo \textbf{Quadro} (definido pelo comando \texttt{\textbackslash begin\{quadro\}}), amplamente utilizado no PPGT para dados qualitativos, conceituais ou comparativos com bordas fechadas:

\begin{quadro}[htbp]
  \caption{Exemplo de quadro estruturado segundo as normas do PPGT}
  \label{quad:exemplo_quadro}
  \centering
  \small
  \begin{tabular}{|p{3cm}|p{6cm}|p{4.5cm}|}
    \hline
    \textbf{DIMENSÃO} & \textbf{DESCRIÇÃO DO CONCEITO} & \textbf{APLICAÇÃO NO SETOR} \\
    \hline
    Planejamento & Definição estratégica e alocação de recursos em infraestrutura. & Priorização de investimentos em corredores de transporte. \\
    \hline
    Operação & Gerenciamento dinâmico e controle de tráfego em tempo real. & Melhoria da fluidez e segurança operacional de vias. \\
    \hline
  \end{tabular}
  \fonte{Elaborado pelo autor (2026).}
\end{quadro}

A seguir é apresentado um exemplo ilustrativo de inserção de \textbf{Figura}:

\begin{figure}[htbp]
  \centering
  % Para inserir uma imagem real, descomente a linha abaixo e aponte para o arquivo correspondente:
  % \includegraphics[width=0.8\textwidth]{imagens/fluxograma_pesquisa.png}
  \framebox[0.7\textwidth]{\rule{0pt}{4cm}\centering\small [Espaço reservado para diagrama ou imagem ilustrativa]}
  \caption{Exemplo de representação gráfica ou fluxograma metodológico}
  \label{fig:exemplo_figura}
  \fonte{Elaborado pelo autor (2026).}
\end{figure}

\section{ESTRUTURA DA DISSERTAÇÃO}
\label{sec:estrutura}

Apresente aqui a organização dos capítulos do trabalho. Por exemplo:
O presente trabalho está estruturado em cinco capítulos, além dos elementos pós-textuais. O Capítulo~\ref{capitulo2} apresenta a revisão bibliográfica e a fundamentação teórica. O Capítulo~\ref{capitulo3} descreve os materiais e métodos empregados. O Capítulo~\ref{capitulo4} detalha a aplicação empírica e a análise dos resultados obtidos. Por fim, o Capítulo~\ref{capitulo5} sintetiza as conclusões, limitações do estudo e sugestões para trabalhos futuros.

  \chapter{REVISÃO BIBLIOGRÁFICA}
\label{capitulo2}

Apresente neste capítulo o estado da arte e a fundamentação teórica que subsidiam a pesquisa. A revisão de literatura deve sintetizar criticamente os principais autores, modelos, teorias e evidências empíricas já estabelecidos na literatura nacional e internacional.

\section{FUNDAMENTAÇÃO TEÓRICA E DIRETRIZES DE CITAÇÃO}
\label{sec:fundamentacao}

O modelo do PPGT adota o sistema de chamada \textbf{Autor-Data}, totalmente gerenciado pelo pacote \texttt{biblatex} integrado ao motor \texttt{biber}. As citações no texto devem seguir estritamente duas modalidades de chamada:

\begin{enumerate}
  \item \textbf{Citação Textual (Narrativa / Indireta integrada ao texto):} Deve ser realizada com o comando \texttt{\textbackslash citet\{chave\}}. Neste formato, o nome do autor é integrado sintaticamente à oração com iniciais em maiúsculo e apenas o ano é posicionado entre parênteses.
  \begin{itemize}
    \item Exemplo com um autor: Segundo \citet{davis2004}, os acidentes de trânsito exigem abordagens desagregadas de análise.
    \item Exemplo com dois autores (unidos automaticamente por \&): Conforme demonstrado por \citet{brackstone1999}, os modelos de perseguição veicular representam aspectos dinâmicos de tráfego.
    \item Exemplo com três ou mais autores (com \textit{et al.} automático em itálico): De acordo com \citet{abdulsalam2015effect}, a automação de processos fiscalizatórios impacta diretamente os indicadores observados.
  \end{itemize}

  \item \textbf{Citação Entre Parênteses (Parentética ao final do parágrafo ou oração):} Deve ser realizada com o comando \texttt{\textbackslash citep\{chave\}}. O sobrenome do autor é automaticamente convertido para \textbf{LETRAS MAIÚSCULAS} dentro dos parênteses.
  \begin{itemize}
    \item Exemplo com um autor: A modelagem matemática de eventos viários requer cuidados metodológicos especiais \citep{davis2004}.
    \item Exemplo com dois autores: O comportamento do motorista no seguimento veicular depende de múltiplos parâmetros perceptuais \citep{brackstone1999}.
    \item Exemplo com três ou mais autores: Estudos empíricos evidenciam o papel de intervenções viárias estruturadas \citep{abdulsalam2015effect}.
  \end{itemize}
\end{enumerate}

\section{USO DE SIGLAS E SÍMBOLOS MATEMÁTICOS}
\label{sec:siglas_simbolos}

Para incluir automaticamente termos na \textbf{Lista de Símbolos, Nomenclatura e Abreviações}, utilize os comandos semânticos disponíveis no template:
\begin{itemize}
  \item Para siglas: \texttt{\textbackslash sigla\{PPGT\}\{Programa de Pós-Graduação em Transportes\}}
  \sigla{PPGT}{Programa de Pós-Graduação em Transportes}
  \sigla{UnB}{Universidade de Brasília}
  \item Para símbolos matemáticos: \texttt{\textbackslash simbolo\{\$\textbackslash epsilon\$\}\{Deformação específica\}}
  \simbolo{$\epsilon$}{Deformação específica}
  \simbolo{$\mu$}{Média amostral ou coeficiente de atrito}
\end{itemize}

\section{TÓPICOS CONCLUSIVOS DA REVISÃO}
\label{sec:conclusivos_revisao}

Sintetize os principais consensos, contradições e lacunas identificados na literatura, demonstrando claramente como a metodologia a ser detalhada no Capítulo~\ref{capitulo3} preenche essas oportunidades de investigação.

  \chapter{MATERIAIS E MÉTODOS}
\label{capitulo3}

Descreva neste capítulo os materiais, ferramentas, bases de dados e procedimentos metodológicos utilizados para conduzir a pesquisa, garantindo rigor científico e reprodutibilidade dos experimentos ou coletas de dados.

\section{MATERIAIS E BASES DE DADOS}
\label{sec:materiais}

Descreva os equipamentos, softwares, bases de dados secundárias e insumos empregados. Indique detalhadamente fontes, períodos temporais de coleta e características técnicas das ferramentas.

\section{MÉTODOS E PROCEDIMENTOS}
\label{sec:metodos}

Apresente as etapas do método de pesquisa, as técnicas de amostragem, modelagem matemática ou simulação adotadas.

\subsection{Padronização de Tabelas no PPGT}
\label{subsec:tabelas}

Diferente dos \textit{Quadros} (que possuem bordas laterais fechadas e abrigam dados predominantemente textuais), as \textbf{Tabelas} são destinadas a dados numéricos e seguem a norma de apresentação tabular do IBGE e do PPGT:
\begin{itemize}
  \item Bordas laterais abertas (sem linhas verticais nas extremidades esquerda e direita);
  \item Linhas horizontais no topo, abaixo dos cabeçalhos das colunas e na base inferior da tabela;
  \item Cabeçalhos de colunas em negrito;
  \item Indicação da fonte na base inferior com o comando \texttt{\textbackslash fonte\{...\}} quando extraída ou adaptada de outros trabalhos (dispensável quando elaborada pelo próprio autor).
\end{itemize}

A Tabela~\ref{tab:exemplo_tabela} ilustra a formatação correta de uma tabela segundo o padrão exigido:

\begin{table}[htbp]
  \centering
  \caption{Exemplo de tabela numérica formatada segundo as normas do PPGT}
  \label{tab:exemplo_tabela}
  \small
  \begin{tabular}{lrr}
    \hline
    \textbf{Categoria de Análise} & \textbf{Frequência Absoluta} & \textbf{Percentual (\%)} \\
    \hline
    Classe A & 142 & 45,8 \\
    Classe B & 98 & 31,6 \\
    Classe C & 70 & 22,6 \\
    \hline
    \textbf{Total} & \textbf{310} & \textbf{100,0} \\
    \hline
  \end{tabular}
  % Se a tabela for elaborada integralmente pelo próprio autor, remova ou comente a linha \fonte abaixo:
  \fonte{adaptado de \citet{dnit2006}}
\end{table}

\subsection{Regra Específica do PPGT para o 3º Nível de Subitens}
\label{subsec:terceiro_nivel}

Uma diretriz exclusiva das normas do PPGT é o tratamento do 3º nível de subseção (\texttt{\textbackslash subsubsection}):
\begin{itemize}
  \item No texto, o título do 3º nível aparece em formato itálico com identificador alfabético: \textbf{(a)}, \textbf{(b)}, \textbf{(c)}... em vez da numeração decimal clássica (ex.: 3.2.1);
  \item No Sumário (Índice), os títulos de 3º nível não são exibidos (a profundidade é limitada a 2 níveis, conforme exigido institucionalmente).
\end{itemize}

\subsubsection{Primeira Etapa do Procedimento Analítico}

Exemplo de texto pertencente ao primeiro subitem do 3º nível, identificado automaticamente pelo prefixo (a).

\subsubsection{Segunda Etapa do Procedimento Analítico}

Exemplo de texto pertencente ao segundo subitem do 3º nível, identificado automaticamente pelo prefixo (b).

  \chapter{ANÁLISE DOS RESULTADOS}
\label{capitulo4}

Apresente neste capítulo os resultados obtidos na pesquisa de maneira objetiva, clara e lógica. Os dados apresentados devem ser interpretados e confrontados com a literatura científica revisada no Capítulo~\ref{capitulo2}.

\section{APRESENTAÇÃO DOS DADOS E RESULTADOS}
\label{sec:apresentacao_resultados}

Apresente os resultados preliminares ou agregados decorrentes da aplicação do método descrito anteriormente. Utilize gráficos, quadros e tabelas para organizar as informações de maneira didática e concisa.

\section{ANÁLISE ESTATÍSTICA E DISCUSSÃO DOS ACHADOS}
\label{sec:analise_discussao}

Discuta o significado dos resultados à luz das hipóteses formuladas no Capítulo~\ref{capitulo1}:
\begin{itemize}
  \item Avalie se as hipóteses iniciais foram confirmadas ou refutadas pelos dados empíricos;
  \item Destaque relações estatisticamente significativas ou comportamentos anômalos encontrados;
  \item Compare os resultados alcançados com os estudos prévios de referência citados na literatura.
\end{itemize}

  \chapter{CONCLUSÕES}
\label{conclusao}
\label{capitulo5}

Sintetize neste capítulo as principais conclusões alcançadas pela pesquisa, respondendo aos objetivos propostos e destacando as contribuições científicas e práticas para a área de Transportes.

\section{LIMITAÇÕES DO TRABALHO}
\label{sec:limitacoes}

Apresente com transparência científica as limitações metodológicas, operacionais, geográficas ou temporais enfrentadas ao longo da condução da pesquisa (por exemplo: tamanho da amostra, resolução espacial dos dados ou restrições de modelos adotados).

\section{RECOMENDAÇÕES PARA TRABALHOS FUTUROS}
\label{sec:trabalhos_futuros}

Aponte direções claras e promissoras para investigações futuras que possam complementar, estender ou aprofundar os achados obtidos nesta dissertação ou tese.

\end{verbatim}

Se você precisar alterar o nome de um capítulo, ou adicionar um capítulo inteiramente novo à sua dissertação, siga estes dois passos:

\begin{enumerate}
    \item Crie o novo arquivo (ou renomeie o existente) diretamente dentro da pasta \texttt{capitulos/}. 
    
    \begin{atencao}
    Jamais utilize espaços, acentos ou caracteres especiais nos nomes dos arquivos (ex: prefira \texttt{capitulo6-discussao.tex} em vez de \texttt{capítulo 6 - discussão.tex}).
    \end{atencao}
    
    \item Atualize o bloco de comandos \texttt{\textbackslash include\{...\}} no \texttt{principal.tex} para refletir o nome exato do novo arquivo, inserindo-o na ordem lógica desejada.
\end{enumerate}

\subsection{Resumo ou Epígrafe do Capítulo (\texttt{\textbackslash resumodocapitulo})}

É possível na redação acadêmica incluir uma breve síntese introdutória (ou epígrafe) logo após o título do capítulo, antes do primeiro parágrafo. Para isso, o template disponibiliza o comando:

{\small
\begin{verbatim}
\resumodocapitulo{Breve síntese ou contexto dos tópicos abordados neste capítulo.}
\end{verbatim}
}

Esse comando posiciona o texto alinhado à direita em formato itálico com largura proporcional de 50\% da página, com espaçamento vertical ajustado. Vale ressaltar que o manual normativo do PPGT não prevê resumo no início de capítulos (trata-se de um recurso opcional mantido por compatibilidade com a FT/UnB). Por essa razão, o comando vem comentado por padrão no arquivo \texttt{capitulos/capitulo1-introducao.tex}, podendo ser ativado caso o seu orientador aprove essa formatação.

\section{Ambientes de Compilação: Onde escrever e compilar seu código}

Assim como o Word precisa do aplicativo instalado no computador para abrir arquivos \texttt{.docx}, o código \LaTeX \ precisa de um ``motor'' (compilador) e de uma interface (editor) para ler seus arquivos \texttt{.tex} e gerar o PDF final. Existem duas abordagens principais para montar o seu ambiente de trabalho:

\subsection{A Rota da Nuvem: Overleaf}

Para a esmagadora maioria dos iniciantes, o \textbf{Overleaf} (\url{www.overleaf.com}) é o caminho com menor atrito. Trata-se de uma plataforma online que processa tudo nos servidores deles. Não requer instalação local, já possui todos os pacotes pré-carregados e salva seu progresso na nuvem.

\subsubsection{Passo a Passo para Carregar o Projeto (.zip) no Overleaf}

\begin{enumerate}
    \item \textbf{Obter o arquivo .zip:} Baixe o pacote compactado \texttt{.zip} contendo todos os arquivos e pastas do template do PPGT.
    \item \textbf{Acessar o Overleaf:} Acesse \url{https://www.overleaf.com/} e faça login em sua conta (ou registre-se gratuitamente).
    \item \textbf{Carregar o Projeto:} No painel inicial (\textit{Dashboard}), clique no botão \textbf{New Project} (Novo Projeto) $\rightarrow$ \textbf{Upload Project} (Carregar Projeto).
    \item \textbf{Arrastar o arquivo:} Arraste e solte o arquivo \texttt{.zip} na janela do navegador (o Overleaf descompactará toda a árvore de diretórios automaticamente).
    \item \textbf{Verificar Configurações do Menu:} Clique no botão \textbf{Menu} no canto superior esquerdo e confirme:
    \begin{itemize}
        \item \textbf{Compiler:} \texttt{pdfLaTeX};
        \item \textbf{TeX Live version:} \texttt{2024} (ou a mais recente);
        \item \textbf{Main document:} \texttt{principal.tex}.
    \end{itemize}
    \item \textbf{Compilar:} Abra o arquivo \texttt{principal.tex} e pressione o atalho \texttt{Ctrl + Enter} (ou clique no botão verde \textbf{Recompile}). O PDF da sua dissertação será renderizado no painel à direita.
    \item \textbf{Edição Modular:} Edite seus dados em \texttt{principal.tex} e redija os capítulos dentro da pasta \texttt{capitulos/}. Todas as imagens devem ser inseridas na pasta \texttt{imagens/}.
\end{enumerate}

\begin{atencao}[Atenção aos arquivos auxiliares no Overleaf]
Se desejar limpar a árvore de arquivos no Overleaf, você pode apagar o arquivo \texttt{Makefile} (que é utilizado apenas em terminais locais de Linux/macOS). \textbf{Contudo, jamais apague o arquivo \texttt{.latexmkrc}}, pois o Overleaf o utiliza internamente para processar e gerar a Lista de Símbolos, Nomenclatura e Siglas (\texttt{nomencl}) de forma automática.
\end{atencao}

\subsection{A Rota Local: MiKTeX, Perl e IDEs (Windows)}

Para pesquisadores que preferem trabalhar \textit{offline}, precisam compilar documentos imensos mais rapidamente ou desejam usar o controle de versão do GitHub de forma nativa, é necessário preparar o ambiente local no Windows. Isso exige três ferramentas trabalhando em conjunto:

\begin{itemize}
    \item \textbf{O Motor Tipográfico (MiKTeX):} O núcleo do sistema. O MiKTeX é a distribuição \LaTeX \ padrão para Windows. Ele traduz os comandos matemáticos e textuais para a tela. Sua maior vantagem é a instalação de pacotes sob demanda (\textit{on-the-fly}): se o seu código pedir um pacote novo, o MiKTeX o baixa da internet automaticamente durante a compilação.
    \item \textbf{A Dependência Oculta (Perl):} O \LaTeX \ depende de scripts externos para gerenciar tarefas pesadas, como compilações contínuas (\texttt{latexmk}), processamento avançado de bibliografias e índices. No Windows, o motor não consegue rodar essas ferramentas nativamente. É estritamente necessário instalar um interpretador, como o \textbf{Strawberry Perl}, para evitar que a compilação trave de forma silenciosa e inexplicável.
    \item \textbf{O Editor de Código (IDE):} Você precisa de um programa desenhado para escrever código confortavelmente. O \textbf{Visual Studio Code (VS Code)} é uma das opções mais poderosas da atualidade; instalando a extensão \textbf{LaTeX Workshop}, você transforma o VS Code em uma máquina de escrever acadêmica com autocompletar de comandos, realce de sintaxe e pré-visualização sincronizada do PDF. Outra alternativa robusta e pronta para uso é o \textbf{TeXstudio}, um software independente inteiramente dedicado ao LaTeX.
\end{itemize}

Após instalar essas três ferramentas, você estará pronto para compilar seus documentos localmente no Windows. Lembre-se de sempre testar a compilação com o arquivo \texttt{principal.tex} para garantir que todos os pacotes e scripts estão funcionando corretamente.

\begin{atencao}[Regra de Ouro no VS Code: Clicar no Documento antes de Compilar]
Ao utilizar o \textbf{Visual Studio Code} com a extensão \textbf{LaTeX Workshop}, há um detalhe operacional fundamental que confunde muitos iniciantes: \textbf{o cursor do teclado/mouse precisa estar ativo dentro do documento \texttt{.tex}} antes de acionar a compilação (seja pelo atalho \texttt{Ctrl+Alt+B}, pelo botão de \textit{Build LaTeX project} na aba lateral do LaTeX ou no menu da barra de status).

Se o foco da janela estiver fora do editor de texto, por exemplo, com o cursor dentro do terminal integrado do VS Code, na árvore de pastas do explorador, no painel de saída (\textit{Output}) ou em um arquivo não-\LaTeX \ (como uma aba de PDF ou arquivo de configurações), a extensão não consegue identificar em qual projeto você está trabalhando e \textbf{o documento simplesmente não compila} (ou acusa que não encontrou o arquivo raiz). Se você clicar para compilar e o processo não iniciar, a solução é simples: \textbf{clique com o mouse dentro de qualquer linha de texto do seu arquivo \texttt{principal.tex} e tente novamente}.
\end{atencao}

\begin{atencao}[Bloqueio de Arquivo no Windows (Adobe Acrobat Reader)]
Caso vá compilar localmente via terminal ou VS Code e o compilador acusar o erro:
\begin{verbatim}
! I can't write on file 'principal.pdf'.
\end{verbatim}
certifique-se de \textbf{fechar a aba do PDF no Adobe Acrobat Reader}. No Windows, o Acrobat bloqueia o arquivo exclusivamente para leitura/gravação enquanto ele estiver aberto na tela.

No visualizador integrado do VS Code (\textit{LaTeX Workshop}) ou em navegadores web (Edge, Chrome, Firefox), esse bloqueio não ocorre e o PDF é atualizado dinamicamente a cada nova compilação.
\end{atencao}

\begin{atencao}[E se o Perl não funcionar ou o \texttt{latexmk} der erro?]
O utilitário \texttt{latexmk} é, por natureza, um script escrito na linguagem Perl. Se o seu computador não tiver o interpretador Perl instalado (ou se o Windows acusar que \textit{``'perl' não é reconhecido como um comando interno''}), o \texttt{latexmk} não conseguirá rodar. Para resolver isso sem estresse, existem duas rotas:

\begin{enumerate}
    \item \textbf{Rota 1: Instalar o Strawberry Perl (Recomendado):} Acesse \url{https://strawberryperl.com/}, baixe o instalador padrão para Windows e reinicie seu terminal ou VS Code. O comando \texttt{latexmk} passará a funcionar de forma totalmente transparente e automática.
    \item \textbf{Rota 2: Compilação Manual Direta (100\% Independente de Perl):} Os programas principais de diagramação do MiKTeX (\texttt{pdflatex.exe}, \texttt{biber.exe} e \texttt{makeindex.exe}) são binários nativos de Windows e \textbf{não necessitam de Perl}. Você pode executar a compilação completa rodando os seguintes comandos no terminal:
\begin{verbatim}
pdflatex -interaction=nonstopmode principal.tex
biber principal
makeindex principal.nlo -s nomencl.ist -o principal.nls
pdflatex -interaction=nonstopmode principal.tex
pdflatex -interaction=nonstopmode principal.tex
\end{verbatim}
    Dica prática: você pode colar essas cinco linhas em um arquivo de texto chamado \texttt{compilar.bat} na raiz da sua dissertação; assim, basta dar um duplo clique nele para recompilar todo o trabalho sem precisar de Perl.
\end{enumerate}
\end{atencao}

\section{Tópicos Avançados: Controle de Versão com GitHub}

Para quem está acostumado exclusivamente com o Microsoft Word, o ciclo de vida de um documento costuma terminar em uma pasta caótica cheia de arquivos nomeados como \textit{dissertacao\_final.docx}, \textit{dissertacao\_final\_revisada.docx}, e \textit{dissertacao\_final\_AGORAVAI.docx}. Além disso, o recurso de ``Controlar Alterações'' (Track Changes) do Word pode tornar o texto visualmente poluído e confuso quando há múltiplas revisões do orientador.

Como o \LaTeX \ baseia-se em arquivos de texto puro (arquivos \texttt{.tex}), ele abre as portas para ferramentas profissionais de engenharia de software, sendo o \textbf{Git} (e a plataforma \textbf{GitHub}) a mais famosa delas. 

Em termos simples, o GitHub atua como uma ``máquina do tempo'' e um diário de bordo superpoderoso para a sua dissertação:

\begin{itemize}
    \item \textbf{O Fim das Cópias Múltiplas:} Você trabalha em um único arquivo (ex: \texttt{principal.tex}). O Git salva o histórico das versões nos bastidores, invisível para você no dia a dia, mantendo sua pasta de trabalho completamente limpa.
    \item \textbf{Fotografias do Progresso (Commits):} Em vez de apenas apertar ``Salvar'', no Git você faz um \textit{commit} — que é como tirar uma fotografia exata de como seu documento estava naquele segundo, acompanhada de um bilhete explicando o que mudou (ex: \textit{``Adicionada a tabela de resultados no capítulo 3''}). 
    \item \textbf{Controle de Alterações Cirúrgico:} Se você apagou um parágrafo importante há três meses e precisa dele de volta, o GitHub permite revisitar qualquer ``fotografia'' do passado. Ele mostra exatamente qual linha foi adicionada (em verde) e qual foi apagada (em vermelho), sem poluir a leitura do seu documento atual.
    \item \textbf{Segurança contra Desastres:} O GitHub hospeda o histórico do seu trabalho na nuvem. Se o seu computador quebrar ou for roubado, todo o seu texto, imagens e bibliografias estão seguros e podem ser baixados em outro computador em segundos.
\end{itemize}

Embora exija o aprendizado de alguns conceitos novos (como \textit{commit}, \textit{push} e \textit{pull}), usar o GitHub transforma a ansiedade de perder dados em uma tranquilidade absoluta, permitindo que você escreva e reescreva sua pesquisa com total liberdade e rastreabilidade.

\section{Apêndice A: Comportamento das Nomenclaturas e Siglas}

A geração das listas de siglas e símbolos costuma ser um ponto de grande frustração na reta final da escrita de uma dissertação. No fluxo tradicional, o pesquisador anota manualmente cada nova sigla em um arquivo separado e depois luta para colocá-las em estrita ordem alfabética, revisando se não esqueceu nenhuma.

No template do PPGT, esse processo é inteiramente abstraído. O comportamento das nomenclaturas segue uma lógica declarativa e descentralizada:

\begin{itemize}
    \item \textbf{Declaração Contextual:} Você não precisa manter uma lista separada. Quando utilizar uma sigla pela primeira vez no meio de qualquer capítulo, declare-a imediatamente usando o comando \texttt{\textbackslash sigla\{IBGE\}\{Instituto Brasileiro de Geografia e Estatística\}}. 
    \item \textbf{Ordenação Alfabética Autônoma:} Não se preocupe com a ordem em que as siglas ou símbolos aparecem no seu texto. Durante o processo de compilação, o motor do \LaTeX \ varre todos os capítulos, coleta cada comando \texttt{\textbackslash sigla} ou \texttt{\textbackslash simbolo} declarado, remove duplicatas automaticamente e processa a lista pré-textual em rigorosa ordem alfabética.
    \item \textbf{Diferenciação Semântica:} Mantenha a organização lógica do seu código utilizando \\ \texttt{\textbackslash sigla\{...\}\{...\}} para instituições e acrônimos, e \texttt{\textbackslash simbolo\{...\}\{...\}} para variáveis matemáticas e grandezas físicas (ex: \texttt{\textbackslash simbolo\{\$v\_m\$\}\{Velocidade média\}}). O sistema separará e construirá a Lista de Abreviaturas e Siglas e a Lista de Símbolos nas páginas corretas.
\end{itemize}

\begin{atencao}[A Compilação das Listas]
Para que a ordenação alfabética seja gerada corretamente, ferramentas auxiliares como o \textit{MakeIndex} entram em ação nos bastidores. Se você estiver usando o Overleaf, isso ocorre de forma invisível. Caso compile localmente, lembre-se de que pode ser necessário compilar o documento duas vezes para que a lista reflita as últimas siglas adicionadas.
\end{atencao}

\section{Apêndice B: Por que não utilizar o pacote \texttt{abntex2}?}

Uma das primeiras inclinações ao modernizar um documento acadêmico brasileiro é importar o pacote \texttt{abntex2} para gerenciar a formatação. Neste projeto, essa prática foi terminantemente descartada por razões arquiteturais e normativas:

\begin{itemize}
    \item \textbf{Divergência Normativa Explícita:} O manual de referências do programa é taxativo ao afirmar que estas normas não seguem o padrão definido pela Associação Brasileira de Normas Técnicas (ABNT). A simples inclusão do \texttt{abntex2} envenenaria o documento com padrões estruturais proibidos pela secretaria do programa, que exige formatações singulares como a ausência de ponto final após o ano e recuo de 1,25 cm nas referências.
    \item \textbf{Legado e Controle de Código:} Mesmo se adotássemos o \texttt{abntex2}, teríamos que modificar profundamente o seu comportamento padrão para forçar o enquadramento nas regras do PPGT. Alterar o núcleo de um pacote externo é tecnicamente custoso e extremamente frágil, pois qualquer atualização futura do \texttt{abntex2} poderia sobrescrever nossas customizações e quebrar a formatação da noite para o dia. Manter nossa própria camada de estilo \\ (\texttt{formatacao/ppgt-unb-format.sty} com BibLaTeX e Biber) isola o documento dessas instabilidades e garante controle total sobre regras críticas.
\end{itemize}

\end{document}